
The special Cosserat rod theory models a directed curve embedded in space
with material points labeled by a coordinate \(\xi\in\mathbb{R}\).
At each point \(\xi\) the current configuration is
identified with a vector \(\bm{x}(\xi)\in\mathbb{R}^3\) describing the
position of the %its 
line of centroids of the directed curve, and two orthogonal directors
\(\left\{\frameBasis_\alpha(\xi)\right\}\) defining a rigid plane called
a cross-section.
A third director \(\frameBasis(\xi)\) is introduced
which is perpendicular to the section,
so that the three directors form %forming
a right-handed orthonormal triad of vectors
\(\{\frameBasis, \frameBasis_\alpha\}\).
Following the approach by \cite{simo1985finite}, this triad 
is identified
abstractly by a single variable \(\FrameRotation(\xi)\) representing the
constrained map from the directors of the reference configuration 
\(\left\{\FrameBasis, \FrameBasis_\alpha\right\}\) to those
in the current configuration \(\{\frameBasis, \frameBasis_\alpha\}\).
The constraints on this map are identical to those of \eqref{eq:def-SO3}
indicating \(\FrameRotation\) can be identified with
the special orthogonal group \(\mathrm{SO}(3)\), and consequently the
problem is without a vectorial configuration manifold. 
%As indicated in
%Section~\ref{sec:rotations}, \(\FrameRotation\) presents algebraically as a
%tensor in the developments that follow, but its computational
%representation is implemented in terms of quaternions. 
The body
coordinate \(\reflenx \in[0, L]\) is taken as the reference arc-length and
the notation
\((~\cdot~)^{\prime} \triangleq \mathrm{d}(~\cdot~) / \mathrm{d} \reflenx\) is adopted. For
prescribed configuration boundary values \((\tilde{\bm{x}}, \tilde{\FrameRotation})\) at
\(\reflenx \in \partial\mathcal{B}\), the configuration manifold has the representation:
\begin{equation*}
\mathscr{C} \triangleq \left\{\,
\chi \triangleq (\bm{x}, \FrameRotation) \colon [0, L] \rightarrow \mathbb{R}^3 \times \operatorname{SO}(3)
\bigm|
\chi = (\tilde{\bm{x}}, \tilde{\FrameRotation})
\text{ on } \partial\mathcal{B}
\,\right\}.
\end{equation*}
The tangent space associated with a configuration \(\chi\) is represented by:
\begin{equation*}
T_\chi \mathscr{C} \triangleq \left\{
(\bm{\eta}_x, \bm{\eta}_{\scriptscriptstyle{\Lambda}})\colon [0, L] \rightarrow \mathbb{R}^3 \times T_{\scriptscriptstyle{\Lambda}}^r \mathrm{SO}(3)
\bigm|
(\bm{\eta}_x, \bm{\eta}_{\scriptscriptstyle{\Lambda}})=(\mathbf{0}, \mathbf{0})
\text{ on } \partial\mathcal{B}
\,\right\}
\end{equation*}
where the right representation \(T_{\scriptscriptstyle{\Lambda}}^r \mathrm{SO}(3)\) of the tangent space on \(\mathrm{SO}(3)\) is used \eqref{eq:rTSO3}.
Leveraging the inner product on \(\mathrm{SO}(3)\)
in \eqref{eq:def-SO3-inner}, the associated inner product on
\(\mathscr{C}\) is given by
\begin{equation} 
\begin{aligned}
\left\langle \bm{\eta}, \bm{u}\right\rangle 
%& \triangleq \int  \bm{\eta}_x \cdot \bm{u}_x + \frac{1}{2} \operatorname{tr}\left[(\bm{\eta}_{\scriptscriptstyle{\Lambda}}^{\times} \FrameRotation)^{\mathrm{t}}(\bm{u}_{\scriptscriptstyle{\Lambda}}^{\times} \FrameRotation)\right]  \,d\xi \\
& =\int  \bm{\eta}_x \cdot \bm{u}_x + \bm{\eta}_{\scriptscriptstyle{\Lambda}} \cdot \bm{u}_{\scriptscriptstyle{\Lambda}} \, \mathrm{d}\reflenx
\end{aligned}
\label{eq:rod-inner}
\end{equation}
for variations
\(\bm{\eta}=(\bm{\eta}_x, \bm{\eta}_{\scriptscriptstyle{\Lambda}})\)
and
\(\bm{u}=(\bm{u}_x, \bm{u}_{\scriptscriptstyle{\Lambda}})\)
in \(T_\chi \mathscr{C}\).
The kinematic transformations $\bm{A}$ and $\bm{A}_*$ are:
\begin{equation}
\bm{A} = \operatorname{diag}(\FrameRotation^{\mathrm{t}}, \FrameRotation^{\mathrm{t}})
\qquad\text{ and }\qquad
\bm{A}_* = \operatorname{diag}(\FrameRotation, \FrameRotation).
\label{eq:cosserat-pull}
\end{equation}


\subsubsection{Equilibrium}

Force \(\shearResult(\xi)\) and moment \(\curveResult(\xi)\)
resultants are defined at each point and the static boundary value
problem enforces the following local momentum balance:
\begin{equation} 
\left\{\begin{array}{r}
\shearResult^{\prime}+{\bm{f}}_n=\mathbf{0} \\
\curveResult^{\prime}+\bm{x}^{\prime} \times \shearResult+{\bm{f}}_m=\mathbf{0} 
\end{array}\right. .
\end{equation}
where \(\bm{f}_n\) and \(\bm{f}_m\) denote given body force and moment loadings, respectively. 
These equations are often motivated by various arguments based on  
finite elasticity theory but 
may also be derived directly from geometric
considerations in connection  %tandem
with a suitable invariance postulate, as is the case in  %was done by
\cite{cosserat1909theorie}.

The classical axial \(N\), shear \(V_\alpha\), torsion \(T\), and moment \(M_\alpha\) resultants 
are represented by \(\shearResult\) and \(\curveResult\) in the director bases:
\begin{equation*}
  \begin{aligned}
    \shearResult
    &= N \frameBasis + V_\alpha \frameBasis_\alpha \\
    &= \FrameRotation\left( N\FrameBasis +  V_\alpha \FrameBasis_\alpha \right) \\
  \end{aligned}
  \qquad \text{ and }\qquad
  \begin{aligned}
    \curveResult
    &= T\, \frameBasis + M_\alpha \frameBasis_\alpha \\
    &= \FrameRotation\left( T\, \FrameBasis + M_\alpha \FrameBasis_\alpha \right) \\
  \end{aligned}
\end{equation*}
where the repeated index \(\alpha\) implies summation over the two
directors %\(\{\frameBasis_{\alpha}\}\)
in the cross section plane.
After collecting the generalized stresses \(\bm{\sigma} = (\shearResult, \curveResult)\)
and body loading \(\bm{f} = (\bm{f}_n, \bm{f}_m)\) 
the strong boundary value problem for static equilibrium is put in the form 
of \eqref{eq:strong} with:
\begin{equation} 
  \mathbb{B} \triangleq -\begin{bmatrix}
    \frac{\mathrm{d}}{\mathrm{d}\xi}\mathbf{1} &  \\
    \bm{x}^{\prime\times} & \frac{\mathrm{d}}{\mathrm{d}\xi}\mathbf{1}
  \end{bmatrix}
  \qquad\text{ so that}\qquad
  \mathbb{B}\bm{\sigma} = -\begin{pmatrix}
    \shearResult^{\prime} \\
    \curveResult^{\prime} + \bm{x}^{\prime}\times \shearResult
  \end{pmatrix}
\label{eq:cosserat-B}
\end{equation}
where $\mathbf{1}$ is the $3 \times 3$ identity and the notation $(~\cdot~)^{\times}$
is introduced in \eqref{eq:hat}.
%
\iffalse
Expanding furnishes on the deformed director basis \(\{\frameBasis_i\}\):
$$
\left.
\begin{aligned}
N^{\prime}\,-\kappa_1 V_1-\kappa_2 V_2+p_t = \dot{\bm{h}}\cdot\frameBasis \\
V_1^{\prime}+\kappa_1 N\,- \tau V_2 + p_1  = \dot{\bm{h}}\cdot\frameBasis_1 \\
V_2^{\prime}+\kappa_2 V_t+ \tau V_1 + f_2  = \dot{\bm{h}}\cdot\frameBasis_2 \\
T^{\prime}\,+\kappa_1 M_2-\kappa_2 M_1 + \gamma_1 V_2-\gamma_2 V_1 + m_t=0 \\
M_1^{\prime} + \kappa_2 T \, - \tau M_2 - \FrameStretch V_2 - \gamma_2 N \, + m_2 = 0 \\
M_2^{\prime} - \kappa_1 T \, + \tau M_1 + \FrameStretch V_1 + \gamma_1 N \, + m_1 = 0 \\
\end{aligned}
\right\}
$$
as presented by \textcite{reissner1973onedimensional}. 
Note however that certain signs differ owing to his peculiar notation for the components of \(\curveResult\). 
That is, from his Equation~3b the components \(M_1\) and \(M_2\) are instead defined by:
$$
\begin{aligned}
\curveResult & = T \frameBasis + \frameBasis \times\left(M_1 \frameBasis_1 + M_2 \frameBasis_2\right) \\
& =T \frameBasis+M_1 \frameBasis_2-M_2 \frameBasis_1
\end{aligned}
$$
\fi 

\subsubsection{Variations}
Applying the inner product \eqref{eq:rod-inner} to the equilibrium \eqref{eq:cosserat-B} furnishes
\begin{equation*}
  \begin{aligned}
  \left<\bm{\eta}, \mathbb{B}\bm{\sigma}\right>
  &= -\int \bm{\eta}_x \cdot \shearResult^{\prime} + \bm{\eta}_{\scriptscriptstyle{\Lambda}}\cdot \left(\curveResult^{\prime} + \bm{x}^{\prime}\times \shearResult\right)\, \mathrm{d} \reflenx \\
  \end{aligned}
\end{equation*}
Integrating by
parts to form \(\left<\mathbb{B}^*\bm{\eta}, \bm{\sigma}\right>\)
gives:
\begin{equation*}
  \begin{aligned}
  \left<\bm{\eta}, \mathbb{B}\bm{\sigma}\right>
  &= -\bm{\eta}_x\cdot\shearResult - \bm{\eta}_{\scriptscriptstyle{\Lambda}}\cdot\curveResult ~ \biggr|_0^L 
  + \int \bm{\eta}_x^{\prime}\cdot\shearResult + \bm{\eta}_{\scriptscriptstyle{\Lambda}}^{\prime}\cdot\curveResult
  - \bm{\eta}_{\scriptscriptstyle{\Lambda}} \cdot(\bm{x}^{\prime}\times \shearResult)\, \mathrm{d} \reflenx
  \end{aligned}
\end{equation*}
Re-writing this in terms of the generalized stress finally reveals the adjoint \(\mathbb{B}^*\):
\begin{equation}
  \begin{aligned}
  \left<\bm{\eta}, \mathbb{B}\bm{\sigma}\right>
  % &= -\bm{\eta}\cdot\bm{\sigma}
  % ~ \biggr|_0^L 
  % + \int \left(\bm{\eta}_x^{\prime} + \bm{x}^{\prime}\times \bm{\eta}_{\scriptscriptstyle{\Lambda}}\right)\cdot \shearResult
  %  + \bm{\eta}_{\scriptscriptstyle{\Lambda}}^{\prime}\cdot\curveResult
  % \,d\xi \\
  &= - \bm{\eta}\cdot\bm{\sigma} ~ \biggr|_0^L + \left<\mathbb{B}^*\bm{\eta}, \bm{\sigma}\right>
  \end{aligned}
  \qquad\text{ where }\qquad 
  \mathbb{B}^* \triangleq \begin{bmatrix}
  \frac{\mathrm{d}}{\mathrm{d} \xi}\mathbf{1} & \bm{x}^{\prime \times}  \\
  \mathbf{0} & \frac{\mathrm{d}}{\mathrm{d} \xi} \mathbf{1}
  \end{bmatrix}.
\label{eq:def-Bstar}
\end{equation}
%
The strain measures arise naturally from the differential geometry of
directed curves embedded in Euclidean space
\citep{cosserat1909theorie,ericksen1957exact}. These
furnish components representing axial \(\epsilon\), shear \(\gamma_{\alpha}\), torsion \(\Twist\),
and curvature \(\kappa_{\alpha}\) strains in the current director basis:
%
\begin{equation*} 
  \begin{aligned}
    \shearStrain &= \bm{x}^\prime - \frameBasis \\
    &= \epsilon \frameBasis + \gamma_\alpha\frameBasis_\alpha
  \end{aligned}
  \qquad\text{ and }\qquad
  \begin{aligned}
    \curveStrain &= \left(\FrameRotation^\prime \FrameRotation^\mathrm{t}\right)^{\vee} - \curveStrain_0 \\
    &= \Twist \frameBasis + \kappa_\alpha \frameBasis_{\alpha}
  \end{aligned}
\end{equation*}
%
where the notation \(\curveStrain^{\times}\) indicates that the quantity \(\FrameRotation^{\prime}\FrameRotation^{\mathrm{t}}\) is skew-symmetric, and the vector \(\curveStrain\) is its axial vector. % is introduced by Equation~(\ref{eq:vee}).
\eqref{eq:cosserat-pull} for $\bm{A}$ furnishes the material stress $\bm{s}$ and strain $\bm{e}$:
% \RepeatEquation{eq:cosserat-strain}
\begin{equation} 
  \bm{s} \triangleq \begin{pmatrix}\FrameRotation^\mathrm{t}\shearResult\\ \FrameRotation^\mathrm{t}\curveResult\end{pmatrix}
  \quad\text{ and }\quad 
  \bm{e} \triangleq \begin{pmatrix}\FrameRotation^\mathrm{t}\shearStrain\\ \FrameRotation^\mathrm{t}\curveStrain\end{pmatrix}.
  % \label{eq:cosserat-strain}
\end{equation}
and one has the related strain vectors in the reference configuration?
\begin{equation*} 
  \begin{aligned}
    \ShearStrain &= \tp{\FrameRotation}\bm{x}^\prime - \tp{\FrameRotation_0}\bm{x}_0^{\prime} \\
    &= \epsilon \FrameBasis + \gamma_\alpha\FrameBasis_\alpha
  \end{aligned}
  \qquad\text{ and }\qquad
  \begin{aligned}
    \CurveStrain^{\times} &= \FrameRotation^\mathrm{t}\FrameRotation^\prime - \FrameRotation_0^\mathrm{t}\FrameRotation_0^\prime \\
    &= \Twist \FrameBasis + \kappa_\alpha \FrameBasis_{\alpha}
  \end{aligned}
\end{equation*}

\iffalse
The Levi-Civita connection on \(\mathscr{C}\) is identified by: 
\begin{equation} 
\nabla_{\bm{u}} \bm{\eta} = \left(\mathbf{0}, \frac{1}{2}(\bm{u}_{\scriptscriptstyle{\Lambda}} \times \bm{\eta}_{\scriptscriptstyle{\Lambda}})^{\times} \FrameRotation\right).
\end{equation}

For arbitrary variations \(\bm{u}, \bm{\eta}\) and \(\bm{\zeta}\) in
\(T_\chi \mathscr{C}\) the well known Koszul formula
\citep{bishop1980tensor} identifies the Levi-Civita connection
associated with a given metric:

\begin{equation} 
\begin{aligned}
2\left\langle\nabla_{\bm{u}} \bm{\eta}, \bm{\zeta}\right\rangle & =\left\langle\left[\bm{u}, \bm{\eta}\right], \bm{\zeta}\right\rangle-\left\langle\left[\bm{u}, \bm{\zeta}\right], \bm{\eta}_\chi\right\rangle-\left\langle\left[\bm{\eta}, \bm{\zeta}\right], \bm{u}_\chi\right\rangle \\
& =\int_0^L(\bm{u}_{\scriptscriptstyle{\Lambda}} \times \bm{\eta}_{\scriptscriptstyle{\Lambda}}) \cdot \bm{\zeta}_{\scriptscriptstyle{\Lambda}} -(\bm{u}_{\scriptscriptstyle{\Lambda}} \times \bm{\zeta}_{\scriptscriptstyle{\Lambda}}) \cdot \bm{\eta}_{\scriptscriptstyle{\Lambda}}-(\bm{\eta}_{\scriptscriptstyle{\Lambda}} \times \bm{\zeta}_{\scriptscriptstyle{\Lambda}}) \cdot \bm{u}_{\scriptscriptstyle{\Lambda}} \, d \xi \\
& =\int_0^L(\bm{u}_{\scriptscriptstyle{\Lambda}} \times \bm{\bm\eta}_{\scriptscriptstyle{\Lambda}}) \cdot \bm{\zeta}_{\scriptscriptstyle{\Lambda}} \,d \xi .
\end{aligned}
\end{equation}
  Therefore the Levi-Civita connection is given by: 
\begin{equation} 
\nabla_{\bm{u}} \bm{\eta} = \left(\mathbf{0}, \frac{1}{2}(\bm{u}_{\scriptscriptstyle{\Lambda}} \times \bm{\eta}_{\scriptscriptstyle{\Lambda}})^{\times} \FrameRotation\right).
\end{equation}
\fi 


